% Chapter 2: Extending Claude — Commands, Skills, Hooks, Agents
\chapter{Extending Claude: Commands, Skills, Hooks, and Agents}
\label{ch:extensions}

Claude Code ships with powerful defaults, but its real leverage comes from
the extension system. Four mechanisms---commands, skills, hooks, and agents---form
a hierarchy of increasing autonomy. Understanding when to use each is the difference
between a tool and a workflow.

\section{Custom Slash Commands}
\label{sec:commands}

\marginnote[Official]{Place in \code{\textasciitilde/.claude/commands/} (global) or \code{.claude/commands/} (project). Invoked as \code{/command-name}.\sourceurl{https://code.claude.com/docs/en/best-practices}}

\official{Commands are Markdown files that expand into prompts when invoked.}
They accept arguments via \code{\$ARGUMENTS} and can reference files with \code{@path}.

\begin{minted}{markdown}
# .claude/commands/review.md
Review the following code for bugs, style violations,
and missing edge cases:

$ARGUMENTS

Focus on: correctness first, then maintainability.
Flag any security concerns.
\end{minted}

\subsection{Design Commands Around Decision Points}

\practitioner{The best commands are not shortcuts---they encode decisions.}
Instead of a command that runs a test, create a command that decides
\emph{what to do next} based on project state.

\begin{minted}{markdown}
# .claude/commands/next.md
Check the current project state:
1. Read CURRENT_WORK.md for context
2. Run `git status` and `git diff --stat`
3. Check for failing tests
4. Recommend the single highest-priority next action
\end{minted}

\practitioner{Chain commands into workflow loops:}

\begin{enumerate}
  \item \code{/next} --- determine what to do
  \item \code{/review} --- verify quality
  \item \code{/ship} --- commit and push
\end{enumerate}

Each command is simple. The workflow emerges from composition.

% -------------------------------------------------------------------
\section{Skills: Domain Knowledge and Workflows}
\label{sec:skills}

\marginnote[Official]{Place in \code{.claude/skills/}. Markdown with YAML frontmatter. Applied automatically or invoked explicitly.\sourceurl{https://code.claude.com/docs/en/best-practices}}

\official{Skills combine domain knowledge with executable workflows.}
Unlike commands (which are prompt templates), skills can:

\begin{itemize}
  \item Be applied automatically when Claude detects a relevant context
  \item Include file globs for automatic triggering
  \item Disable model invocation for side-effect-only workflows
  \item Define multi-step procedures with validation gates
\end{itemize}

\begin{minted}{yaml}
---
name: qa-health
description: Run quality dashboard across all volumes
allowed_tools: [Bash, Read, Glob]
---
# QA Health Check
1. Run `python scripts/audits/health_dashboard.py`
2. Parse output for RED/YELLOW/GREEN status
3. Report findings with actionable recommendations
4. If any RED items, create TODO list for fixes
\end{minted}

\subsection{High-Value Skill Patterns}

\practitioner{Skills that orchestrate multiple scripts produce the highest value.}

\begin{description}
  \item[Quality Dashboard] Run lint + tests + coverage + metrics in one invocation.
    Present traffic-light summary instead of raw output.
  \item[Pre-Release Check] Validate documentation, test coverage, version bumps,
    and changelog entries before creating a release.
  \item[Evidence Verification] Cross-reference claims in documents against
    authoritative sources. Flag unverified assertions.
\end{description}

\margintip{A skill that saves 5 minutes per use and runs 3 times per week
saves 13 hours per year. Measure value by frequency times savings.}

% -------------------------------------------------------------------
\section{Hooks: Deterministic Automation}
\label{sec:hooks}

\marginnote[Official]{Multiple lifecycle events. Configure via \code{/hooks} or in \code{settings.json}. See \cref{app:reference} for the full list.\sourceurl{https://code.claude.com/docs/en/hooks-guide}}

\convergence{Hooks are the enforcement layer. Unlike CLAUDE.md instructions
(which are advisory), hooks \textbf{always execute}.}

\subsection{Available Lifecycle Events}

\begin{fullwidth}
\begin{tabular}{@{}lp{8cm}@{}}
\toprule
\textbf{Event} & \textbf{Use Case} \\
\midrule
\code{PreToolUse} & Validate or transform tool inputs before execution \\
\code{PostToolUse} & React to tool outputs (logging, metrics) \\
\code{PreFileWrite} & Block writes to protected files; enforce formatting \\
\code{PostFileWrite} & Auto-format, lint, or validate written files \\
\code{PreBashRun} & Validate commands before execution; block dangerous ops \\
\code{PostBashRun} & Capture output for logging or metrics \\
\code{PreCommit} & Run test suite; block commits with failing tests \\
\code{PostCommit} & Trigger CI, update dashboards, send notifications \\
\code{UserPromptSubmit} & Transform or validate user input \\
\code{Notification} & Route notifications to Slack, email, or dashboard \\
\bottomrule
\end{tabular}

\vspace{2pt}
{\footnotesize\textit{Selected events shown. See \cref{app:reference} for the complete list
including SessionStart, SessionEnd, WorktreeCreate, and others.}}
\end{fullwidth}

\subsection{Hooks Enforce What CLAUDE.md Can Only Suggest}

\practitioner{The golden rule: if a standard is non-negotiable, make it a hook.
If it is advisory, put it in CLAUDE.md.}

\begin{minted}{json}
{
  "hooks": {
    "PostFileWrite": [{
      "command": "prettier --write \"$FILE\"",
      "description": "Auto-format on every file write"
    }],
    "PreCommit": [{
      "command": "pytest tests/ --tb=short -q",
      "description": "Block commits if tests fail"
    }]
  }
}
\end{minted}

\begin{warnbox}[Hook vs.\ CLAUDE.md]
CLAUDE.md: ``Always run tests before committing.'' \\
Claude may forget, especially in long sessions with degraded context.

Hook: \code{PreCommit -> pytest} \\
Tests run every time, regardless of context state. Cannot be forgotten.
\end{warnbox}

% -------------------------------------------------------------------
\section{Subagents and Agent Teams}
\label{sec:agents}

\subsection{Subagents: Isolated Context Workers}

\marginnote[Official]{Define in \code{.claude/agents/}. Each agent runs with its own context, isolated from the main session.\sourceurl{https://code.claude.com/docs/en/best-practices}}

\official{Subagents are purpose-built Claude instances that run in isolated context.}
They receive a task, execute it with their own tool set, and return results
to the parent session.

Use subagents when:
\begin{itemize}
  \item The task requires heavy context that would pollute your main session
  \item You need an unbiased second opinion (writer/reviewer pattern)
  \item The task can run independently in the background
\end{itemize}

\begin{minted}{markdown}
# .claude/agents/security-reviewer.md
---
name: security-reviewer
description: Review code for security vulnerabilities
allowed_tools: [Read, Grep, Glob]
---
You are a security specialist. Review the provided code for:
1. Injection vulnerabilities (SQL, command, path traversal)
2. Authentication/authorization gaps
3. Sensitive data exposure
4. Dependency vulnerabilities
Report findings with severity ratings (Critical/High/Medium/Low).
\end{minted}

\subsection{Agent Teams: Collaborative Multi-Agent Work}

\marginnote[Official]{Research preview. Enable with environment variable. Multiple agents share a task list and can message each other.\sourceurl{https://code.claude.com/docs/en/agent-teams}}

\official{Agent Teams extend subagents with collaboration: shared task lists,
direct messaging between agents, and team lead orchestration.}

Best suited for:
\begin{itemize}
  \item Research tasks requiring multiple perspectives
  \item Parallel debugging across different system layers
  \item Cross-functional coordination (frontend + backend + database)
\end{itemize}

\subsection{The Delegation Hierarchy}

\practitioner{Match the complexity of your delegation mechanism to the task:}

\begin{enumerate}
  \item \textbf{Commands} --- simple prompt expansion, no isolation
  \item \textbf{Skills} --- multi-step workflows with domain knowledge
  \item \textbf{Subagents} --- isolated context, independent execution
  \item \textbf{Agent Teams} --- collaborative, shared state, multi-agent
\end{enumerate}

% -------------------------------------------------------------------
\section{Visual: The Delegation Hierarchy}

\begin{figure}[H]
\centering
\begin{tikzpicture}[
  box/.style={draw=PrimaryBlue, fill=PrimaryBlue!8, rounded corners=3pt,
    minimum width=3cm, minimum height=0.9cm, font=\small, align=center},
  desc/.style={font=\footnotesize, text=DarkText, align=left},
  arr/.style={->, thick, PrimaryBlue!60},
  >=Stealth
]
  % Boxes arranged vertically
  \node[box] (cmd) at (0,0) {\textbf{Commands}};
  \node[box] (skill) at (0,1.8) {\textbf{Skills}};
  \node[box, fill=PrimaryBlue!12] (sub) at (0,3.6) {\textbf{Subagents}};
  \node[box, fill=PrimaryBlue!16] (team) at (0,5.4) {\textbf{Agent Teams}};

  % Descriptions
  \node[desc, anchor=west] at (2.2,0) {Prompt template, inline context\\
    \textit{Use for: simple shortcuts}};
  \node[desc, anchor=west] at (2.2,1.8) {Multi-step workflow, domain knowledge\\
    \textit{Use for: repeated procedures}};
  \node[desc, anchor=west] at (2.2,3.6) {Isolated context, independent execution\\
    \textit{Use for: research, review, background}};
  \node[desc, anchor=west] at (2.2,5.4) {Shared task list, agent messaging\\
    \textit{Use for: cross-functional collaboration}};

  % Arrows
  \draw[arr] (cmd) -- (skill);
  \draw[arr] (skill) -- (sub);
  \draw[arr] (sub) -- (team);

  % Side labels
  \node[font=\footnotesize\itshape, text=PrimaryBlue!70, rotate=90, anchor=south] at (-2.2,2.7) {Increasing autonomy $\rightarrow$};
  \node[font=\footnotesize\itshape, text=WarningOrange!80, rotate=90, anchor=north] at (-2.8,2.7) {$\leftarrow$ Increasing overhead};
\end{tikzpicture}
\caption{Match delegation mechanism complexity to task requirements.}
\label{fig:delegation-hierarchy}
\end{figure}

\begin{keyinsight}
Most developers jump straight to subagents when a skill would suffice.
The simpler mechanism that solves the problem is the better choice:
less overhead, less context cost, faster execution.
\end{keyinsight}
